\documentclass[10pt,a4paper]{amsart}
\usepackage[margin=19mm]{geometry}
\usepackage{amsmath,amssymb,mathtools}
\usepackage[scr=rsfs,cal=euler]{mathalfa}
\usepackage{xfrac}
\usepackage[unicode,hidelinks,bookmarks=false]{hyperref}
\newcommand{\ie}{i.e.\ }
\newcommand{\cf}{cf.\ }
\newcommand{\resp}{respectively\ }
\newcommand{\Subsection}{\subsection}
\providecommand{\nobreakdash}{\nobreak\hbox{-}}
\setlength{\emergencystretch}{2em}
\multlinegap0pt
\usepackage{bm}
\newtheorem{theorem}{Theorem}
\newcommand{\E}{\mathbb{E}}
\newcommand{\lpar}{\left(}
\newcommand{\rpar}{\right)}
\newcommand{\sumn}{\sum\limits_{i=1}^n}
\title{Sharp Decoupling Inequalities for the Variances and Second Moments of Sums of Dependent Random Variables}
\author{Victor H. de la Pena, Heyuan Yao, Demissie Alemayehu}
\date{Source: arXiv:2512.19063v2 (CC BY 4.0)}
\begin{document}
\maketitle

\noindent\emph{Source and licence.} Sharp Decoupling Inequalities for the Variances and Second Moments of Sums of Dependent Random Variables. Version arXiv:2512.19063v2 (CC BY 4.0), distributed by arXiv under the Creative Commons Attribution 4.0 International licence, http://creativecommons.org/licenses/by/4.0/, as recorded in the arXiv licence metadata for this exact version and shown on the abstract page https://arxiv.org/abs/2512.19063v2. Full licence notice: \texttt{licenses/arxiv-2512.19063-CC-BY-4.0.txt}.\par\smallskip

\noindent\textit{Editorial scope.} Counted unit: the article's theorem decoupling (source0.tex lines 276--283). The article's complete original proof of it is delivered with it (source0.tex lines 284--292, one \texttt{proof} environment of 9 lines). No mathematics is written, repaired or re-derived: the statement and the proof are the article's own lines, in the article's own notation, and no retained line is substituted or re-flowed.  No further source-local context is needed: every cross-reference of every retained line resolves inside the two delivered spans.  Nothing was omitted from any retained span, so the package carries no omission record.  No retained line contains a citation, so the wrapper carries no bibliography and no bibliography entry.  No retained line contains a figure environment, an image-inclusion command or drawing code, so no image is delivered and the package declares no asset dependency.  Editorial formatting only, in addition: an AMS \texttt{amsart} preamble that restates the article's own declarations and macros used by the retained lines, namely the environments \texttt{theorem} on separate counters, together with the standard packages those lines need.  The retained environments print in this document as Theorem 1 in order of appearance, whereas the article prints the same items with the numbering of its own full document.  The complete native source of the article as distributed in the arXiv e-print for this exact version is bundled unchanged as \texttt{sources/191444/source0.tex}.
\begin{theorem}
    \label{thm: complete decoupling}
    For any nonnegative dependent square-integrable  variables $d_1,...d_n$, and the mutually independent  variables $z_1,...,z_n$ such that $z_i \stackrel{\mathcal{L}}{=} d_i$ for all $i=1,...,n$, we have the following sharp inequality
    \begin{equation}
        \label{equ: complete decoupling}
        \frac{1}{2}\E\lpar\sumn z_i\rpar^2 \leq \E \lpar \sumn d_i\rpar^2.
    \end{equation}
\end{theorem}

\begin{proof}
     We notice that $\E (\sumn z_i) ^2 =  \sumn \E z_i^2 + \sum\limits_{1\leq i<j\leq n} \E [z_i z_j] $ and $\E (\sumn d_i) ^2 =  \sumn \E d_i^2 + \sum\limits_{1\leq i<j\leq n} \E [d_i d_j]$. Therefore, since all random variables we consider are nonnegative, we have 
     \[ \E (\sumn d_i) ^2 \geq \sumn \E d_i^2 =\sumn \E z_i^2, \]
    and by Jensen's inequality, 
    \[ \E (\sumn d_i) ^2  \geq (\sumn \E d_i)^2 =(\sumn \E z_i)^2 \geq \sum\limits_{1\leq i<j\leq n} \E [z_i]\E[z_j]  =\sum\limits_{1\leq i<j\leq n}  \E [z_i z_j] . \]
    By adding the above two inequalities, we claim \eqref{equ: complete decoupling}. 
    
    The sharpness of \eqref{equ: complete decoupling} can be seen in the following example. Consider $(d_1,...,d_n) = \boldsymbol{u}_k$ with probability $1/n$, for $k=1,...,n$, where $\boldsymbol{u}_k$ denotes the $n$-dimensional unit vector with the $k$-th element being 1. Therefore, $\sumn d_i \equiv 1$, $\sumn z_i \stackrel{\mathcal{L}}{=} \text{Binomial}(n,1/n)$. Hence $\E (\sumn d_i)^2 =1$ and $\E (\sumn e_i)^2 =2-1/n$. This suggests that the best constant for this lower bound working for all $n$ is $1/2$. 
\end{proof}

\end{document}
